\chapter{Momentum}\label{ch:s1:momentum}

From March to May 2009 the US momentum factor lost 49.4\%, and 52.8\% over the calendar year. It was short the stocks that had fallen most in the crisis, and
those stocks rebounded hardest when the market turned. Momentum's long-run record had looked like one of the best premiums in equities: 7.4\% a year since
1927 at a Sharpe ratio of 0.45. But it earns that record with a skewness of $-3$, and in 1932 it had lost 64.5\% in a year. Scaled each month by its
own recent volatility, the same factor would have lost 16.8\% in 2009 and doubled its long-run Sharpe ratio. This chapter is about momentum in stocks, its
variants, and above all its crashes. The build is \texttt{firm.momstrat}.

\section{Cross-sectional momentum}

\begin{definition}[Cross-sectional momentum]\label{def:s1:momentum:cross}
\emph{Cross-sectional momentum}\index{cross-sectional momentum} is the tendency of stocks with high returns over the past three to twelve months, relative to
other stocks, to keep outperforming over the next months; the standard strategy ranks stocks on their return from twelve months to one month ago (skipping the
most recent month, which reverses) and buys the top and sells the bottom.
\end{definition}

Jegadeesh and Titman documented it: buying past winners and selling past losers earned significant returns over holding periods of three to twelve
months, not explained by systematic risk, and part of the first year's gain faded over the following two. The skipped month is chapter~2's reversal: over one
month, returns reverse, and a momentum signal that includes it trades against itself. The Kenneth French momentum factor is the public benchmark: the
high-minus-low prior-return portfolios among six value-weighted size and prior-return portfolios, rebuilt monthly.

\begin{center}
\small
\begin{tabular}{@{}lcc@{}}
\toprule
French momentum factor, 1927--2026 & raw & scaled\\
\midrule
return a year; volatility & 7.4\%; 16.3\% & 15.4\%; 17.7\%\\
Sharpe ratio & 0.45 & 0.87\\
skewness; excess kurtosis (monthly) & $-3.02$; 27.3 & $-0.40$; 3.1\\
worst month & $-52.6\%$ (8/1932) & $-29.1\%$ (9/1939)\\
maximum drawdown & $-78.4\%$ (6/1932 to 9/1939) & $-42.5\%$ (6/1932 to 5/1933)\\
March to May 2009; calendar 2009 & $-49.4\%$; $-52.8\%$ & $-16.5\%$; $-16.8\%$\\
\bottomrule
\end{tabular}
\end{center}

On the synthetic market, which plants a persistent drift in each stock (Book~7, chapter~5), the 12--1 signal has a monthly rank information coefficient of
0.020, and a decile long--short book (gross exposure one, rebuilt every 21 days) returns 4.1\% a year net of ten basis points per unit traded, with a Sharpe
ratio of 0.39 before costs and 0.31 after and a maximum drawdown of 40\%. It trades 5.7 times its gross a year, one-way.

\section{Residual and industry momentum}

\begin{definition}[Residual momentum, industry momentum]\label{def:s1:momentum:residual}
\emph{Residual momentum}\index{residual momentum} ranks stocks on the sum of their factor-model residual returns over the formation window, divided by the
residuals' standard deviation, instead of on their total returns. \emph{Industry momentum}\index{industry momentum} ranks industries on their past returns
and buys the stocks of the winning industries and sells those of the losing ones.
\end{definition}

A stock's twelve-month return is mostly its factors': a winner in a year when small value stocks rallied is a small value stock. Total-return momentum
therefore carries the factor bets of the past year, and they change every month. Blitz, Huij and Martens found that ranking on residual returns removed these
time-varying exposures and roughly doubled the risk-adjusted profits, which were also steadier and less concentrated in the extreme stocks. Moskowitz and
Grinblatt argued the other way round that industries matter: a strong momentum effect in industry components of returns accounted for much of individual
stock momentum, and \omterm{def:s1:momentum:residual}{industry momentum} strategies were highly profitable on their own.

\begin{center}
\small
\begin{tabular}{@{}lccc@{}}
\toprule
synthetic market, years 3 to 10 & total return & residual & industry\\
\midrule
rank IC, next month & 0.020 & 0.023 & 0.001\\
Sharpe ratio before, after 10 bp & 0.39, 0.31 & 0.55, 0.45 & $-0.04$, $-0.23$\\
return a year after costs; volatility & 4.1\%; 13.4\% & 4.9\%; 11.0\% & $-1.4\%$; 6.4\%\\
maximum drawdown & $-40\%$ & $-33\%$ & $-23\%$\\
\bottomrule
\end{tabular}
\end{center}

The residual version wins for Blitz and co-authors' reason: the same planted drift, less factor noise, a volatility of 11.0\% against 13.4\%. \omterm{def:s1:momentum:residual}{Industry
momentum} earns nothing, because the synthetic market plants none: its industry factors are independent from day to day. That is the synthetic market's
limitation, not evidence against Moskowitz and Grinblatt; it is also a reminder that a signal can only find what its market contains.

\begin{omfigure}
\begin{tikzpicture}
\begin{axis}[width=10cm,height=5.0cm,xlabel={\small year},ylabel={\small cumulative net return (\%)},tick label style={font=\footnotesize},
  xmin=2,xmax=10,grid=major,grid style={gray!25},legend pos=north west,legend style={font=\scriptsize},table/col sep=comma]
\addplot[omDef!60,thick] table[x=year,y=total]{figdata/strategies-1/05-momentum/books.csv};
\addplot[omDef,very thick] table[x=year,y=residual]{figdata/strategies-1/05-momentum/books.csv};
\addplot[omThm,thick] table[x=year,y=industry]{figdata/strategies-1/05-momentum/books.csv};
\legend{total-return momentum,residual momentum,industry momentum}
\end{axis}
\end{tikzpicture}

\omcaption{The three momentum books on the synthetic market, net of ten basis points per unit traded: cumulative returns (gross exposure one). Data:
\texttt{s1\_momentum.book}.}\label{fig:s1:momentum:books}
\end{omfigure}

\section{Intraday and time-of-day momentum}

Momentum also runs at the scale of a day, in the market rather than across stocks. Gao, Han, Li and Zhou found, on S\&P 500 ETF data from 1993 to 2013, that
the market's first half-hour return, measured from the previous close, predicts its last half-hour return, more so on volatile, high-volume, recession and
macroeconomic-news days, and in ten other actively traded ETFs. The explanation they favour is a mix of infrequent portfolio rebalancing and late-informed
trading near the close. It is a futures and ETF strategy on intraday data, which the synthetic daily market cannot test; its strategy file rests on the
published record, and chapter~25 returns to short-horizon futures strategies.

\section{Crashes and their management}

\begin{definition}[Momentum crash]\label{def:s1:momentum:crash}
A \emph{momentum crash}\index{momentum crash} is a short period of very large losses of a momentum strategy, typically when the market rebounds sharply after
a decline, as the past losers the strategy is short (often high-beta stocks beaten down by the decline) rise more than the past winners it holds.
\end{definition}

Daniel and Moskowitz characterised the crashes: infrequent and persistent strings of negative returns, partly forecastable, occurring in panic states,
after market declines and when volatility is high, at the same time as market rebounds. The mechanism is a beta that the strategy acquires without choosing
it. After a long decline, the past losers are the stocks that fell most, which are those with the highest betas; a momentum book is then short high beta,
and a sharp rebound hits it from both sides. The ten worst months of the French factor are six in 1931--1939, April 2009, January 2001, November 2002 and
January 2023. The synthetic market has too few bear markets to show the mechanism (the market fell over the previous two years on only 6.7\% of its days),
which is why this section runs on the French data.

\begin{definition}[Volatility-scaled momentum]\label{def:s1:momentum:scaled}
\emph{Volatility-scaled momentum}\index{volatility-scaled momentum} holds a momentum portfolio with a weight inversely proportional to the strategy's own
recent realised volatility, so that its forecast volatility is constant; the weight is recomputed each period from data before it.
\end{definition}

Barroso and Santa-Clara found that momentum's risk is highly variable and predictable, and that managing it virtually eliminates crashes and nearly doubles
the Sharpe ratio. The chapter's version holds each month a weight of 12\% divided by the annualised volatility of the daily factor over the previous 126
trading days. It lifts the Sharpe ratio from 0.45 to 0.87, the skewness from $-3.02$ to $-0.40$, and the 2009 loss from 52.8\% to 16.8\%. The crashes
come after volatile months: in April 2009, when the raw factor lost 34.4\%, the scaled weight was 0.29, and the scaled loss 10.1\%. Scaling does not help
where the crash comes from calm: in September 1939 the weight was 0.92, and the scaled factor had its worst month ($-29.1\%$). On the synthetic market,
where momentum's volatility varies less, scaling does little: the total-return book's Sharpe ratio stays at 0.52 over the scaled sample and its maximum
drawdown falls from 30\% to 25\%.

\begin{omfigure}
\begin{tikzpicture}
\begin{axis}[width=12.5cm,height=5.4cm,xlabel={\small year},ylabel={\small calendar-year return (\%)},tick label style={font=\footnotesize},
  xmin=1927,xmax=2026,ymin=-70,ymax=150,grid=major,grid style={gray!25},legend pos=north west,legend style={font=\scriptsize},table/col sep=comma,
  x tick label style={/pgf/number format/1000 sep={}}]
\addplot[omThm,thick,mark=*,mark size=0.9pt] table[x=year,y=raw_pct]{figdata/strategies-1/05-momentum/annual.csv};
\addplot[omDef,thick,mark=square*,mark size=0.9pt] table[x=year,y=scaled_pct]{figdata/strategies-1/05-momentum/annual.csv};
\legend{raw,volatility-scaled}
\end{axis}
\end{tikzpicture}

\omcaption{Calendar-year returns of Kenneth French's momentum factor, raw and scaled each month to 12\% by its trailing 126-day volatility, 1928--2025.
Derived from the Kenneth R. French Data Library (Mom, monthly and daily, 202607 CRSP file); the raw series is not
redistributed.}\label{fig:s1:momentum:annual}
\end{omfigure}

The scaled version is not free. It is levered whenever momentum's trailing volatility is below 12\% (much of its mean return of 15.4\%
against the raw factor's 7.4\% comes from that leverage), and a fund running it needs the capacity and financing for it. And the improvement is measured on the same century that suggested it; Daniel and
Moskowitz's dynamic version, which also forecasts momentum's mean from the market's state, reports a similar doubling on several markets and asset classes.

\section{Implementation: turnover, skipping a month, costs}

Momentum is a slow strategy by the standard of this part of the book, and its costs are small per unit of capital: the synthetic decile book trades 5.7 times
its gross a year one-way, which at ten basis points per unit traded costs 1.14\% a year, against 8\% to 38\% for the daily books of chapters~2 and~3. Three
implementation choices matter more than the cost rate. The skipped month removes the reversal, as above. The rebalancing frequency trades off freshness
against turnover (monthly is the convention). And the weighting (equal or by capitalisation, deciles or terciles) moves the book toward small, less liquid
stocks where momentum is stronger and trading dearer; the French factor averages big and small stocks to keep both.

\section{Strategy files}

\begin{strategyfile}{12-1 cross-sectional momentum}\label{strat:s1:momentum:classic}
\sfield{Who pays you, and why} Investors who under-react to news and then chase it; the book is early to the trend and late to its end.
\sfield{Instruments and venues} Liquid stocks, long winners and short losers, or long-only tilts.
\sfield{Signal} Return from twelve months to one month ago, ranked across stocks.
\sfield{Sizing and execution} Deciles or terciles, rebuilt monthly; value or equal weights.
\sfield{Costs} Several times the book traded a year; small per unit of capital.
\sfield{How it dies} Crashes after bear markets when losers rebound (2009, 1932); decay as capital crowds into it (comomentum, Book~7, chapter~28).
\sfield{Horizon, capacity, infrastructure} Months; large capacity; a monthly ranking.
\sfield{Backtest honestly} Point-in-time universes including delisted stocks; the skip month; the crash years in the sample.
\sfield{Sources} Jegadeesh and Titman (1993); French momentum factor, 1927--2026: 7.4\% a year, Sharpe ratio 0.45 before costs, maximum drawdown 78.4\%.
\end{strategyfile}

\begin{strategyfile}{Residual momentum}\label{strat:s1:momentum:residual}
\sfield{Who pays you, and why} As for momentum, with the stock-specific part of the trend isolated.
\sfield{Instruments and venues} As for momentum.
\sfield{Signal} Sum of factor-model residual returns over months $-12$ to $-2$, divided by their standard deviation.
\sfield{Sizing and execution} As for momentum; the book is closer to neutral to the factors by construction.
\sfield{Costs} As for momentum.
\sfield{How it dies} As momentum, with smaller crashes since the book carries less beta.
\sfield{Horizon, capacity, infrastructure} Months; a factor model.
\sfield{Backtest honestly} Residuals from a model estimated only on past data.
\sfield{Sources} Blitz, Huij and Martens (2011): risk-adjusted profits about twice those of total-return momentum.
\end{strategyfile}

\begin{strategyfile}{Industry momentum}\label{strat:s1:momentum:industry}
\sfield{Who pays you, and why} Slow diffusion of industry-wide news.
\sfield{Instruments and venues} Industry baskets or ETFs, or stocks by industry.
\sfield{Signal} Industries' past six- to twelve-month returns.
\sfield{Sizing and execution} Long the top industries, short the bottom; monthly.
\sfield{Costs} Low when traded through industry funds.
\sfield{How it dies} As momentum; concentrated bets on few industries.
\sfield{Horizon, capacity, infrastructure} Months; large capacity.
\sfield{Backtest honestly} Industry classifications as they were at each date.
\sfield{Sources} Moskowitz and Grinblatt (1999).
\end{strategyfile}

\begin{strategyfile}{Intraday momentum in the last half hour}\label{strat:s1:momentum:intraday}
\sfield{Who pays you, and why} End-of-day rebalancers and late-informed traders who trade in the direction of the day's early move.
\sfield{Instruments and venues} Index futures and ETFs.
\sfield{Signal} The first half-hour return from the previous close.
\sfield{Sizing and execution} Position for the last half hour in its direction; flat at the close.
\sfield{Costs} One round trip a day in the most liquid instruments.
\sfield{How it dies} Competition in the closing half hour; changes in who rebalances at the close.
\sfield{Horizon, capacity, infrastructure} Half an hour; intraday data and execution.
\sfield{Backtest honestly} Trade at prices available at the start of the last half hour, not at its average.
\sfield{Sources} Gao, Han, Li and Zhou (2018): significant predictability in S\&P 500 ETF data, 1993--2013.
\end{strategyfile}

\begin{strategyfile}{Volatility-scaled momentum}\label{strat:s1:momentum:scaled}
\sfield{Who pays you, and why} As for momentum; the scaling removes risk the premium does not pay for.
\sfield{Instruments and venues} As for momentum, with leverage available.
\sfield{Signal} The momentum book's weight: a target volatility over its trailing realised volatility.
\sfield{Sizing and execution} Recomputed monthly (or more often) from daily returns before the month.
\sfield{Costs} The momentum book's plus the changes of weight.
\sfield{How it dies} Crashes from calm states (September 1939); leverage constraints in calm years.
\sfield{Horizon, capacity, infrastructure} Months; daily strategy returns.
\sfield{Backtest honestly} Weights from data before each month; financing of the leverage.
\sfield{Sources} Barroso and Santa-Clara (2015); Daniel and Moskowitz (2016); French data: Sharpe ratio 0.45 to 0.87 and a 2009 loss of 16.8\% instead of
52.8\%, 1927--2026, before costs.
\end{strategyfile}

\section{Tutorial: spring 2009}\label{tut:s1:momentum:tut}

\begin{tutorial}
\textbf{Goal.} Build three momentum books on the synthetic market, then measure the 2009 crash in the French factor, raw and scaled. \textbf{End state:} the
two tables and \cref{fig:s1:momentum:annual}.
\begin{tutsteps}
\item \textbf{Signals}: residual and \omterm{def:s1:momentum:residual}{industry momentum}.
\omcode{code/firm/momstrat/firm_momstrat.py}{45}{61}{Residual and industry momentum.}\label{lst:s1:momentum:signals}
\item \textbf{Books}: \texttt{s1\_momentum.book(kind)} and \texttt{summary(kind)} for the three signals.
\item \textbf{Scaling}: a weight from the strategy's own trailing volatility.
\omcode{code/firm/momstrat/firm_momstrat.py}{79}{84}{Volatility scaling.}\label{lst:s1:momentum:scale}
\item \textbf{French data}: \texttt{s1\_fetch\_mom.py} once, then \texttt{french()}; \texttt{fig\_momentum.py}.
\end{tutsteps}
\textbf{What to change next.} Condition the scaling on the market's state (Daniel and Moskowitz); use terciles and value weights like the French factor; add
chapter~2's reversal to the skipped month and measure what it costs.
\end{tutorial}

\section{Build: momentum strategies}\label{bld:s1:momentum:momstrat}

\begin{build}
\bfield{Purpose} The momentum signal family and the crash management every momentum book needs.
\bfield{Interface} \texttt{total\_mom(ret, lookback, skip)}, \texttt{residual\_mom(E, lookback, skip)}, \texttt{industry\_mom(ret, industry, cap, lookback,
skip)}, \texttt{decile\_book(signal, universe, q)}, \texttt{vol\_scale(r, target, window, periods)}, \texttt{bear(mkt, window)}.
\bfield{Rules} Signals and weights from data before each trade; the skip month; scaling from the strategy's own returns.
\bfield{Acceptance tests} \texttt{code/firm/momstrat/tests/}: the skip window by hand; residual and industry signals on planted data; decile weights; the
scaling weight by hand; the market-state flag.
\bfield{Stretch} Dynamic weights from forecasts of the mean and variance (Daniel and Moskowitz); time-series momentum across futures (chapter~19).
\end{build}

\begin{omsources}
\item N.~Jegadeesh and S.~Titman, ``Returns to buying winners and selling losers'', \emph{Journal of Finance} 48(1), 1993.
\item K.~Daniel and T.~J.~Moskowitz, ``Momentum crashes'', \emph{Journal of Financial Economics} 122(2), 2016.
\item P.~Barroso and P.~Santa-Clara, ``Momentum has its moments'', \emph{Journal of Financial Economics} 116(1), 2015.
\item D.~Blitz, J.~Huij and M.~Martens, ``Residual momentum'', \emph{Journal of Empirical Finance} 18(3), 2011.
\item T.~J.~Moskowitz and M.~Grinblatt, ``Do industries explain momentum?'', \emph{Journal of Finance} 54(4), 1999.
\item L.~Gao, Y.~Han, S.~Z.~Li and G.~Zhou, ``Market intraday momentum'', \emph{Journal of Financial Economics} 129(2), 2018.
\item Kenneth R. French Data Library, Momentum Factor (Mom), monthly and daily.
\end{omsources}

\section{Exercises}

\begin{exercise}[$\star$]\label{exo:s1:momentum:1}
The momentum factor's trailing volatility is 40\% a year. What weight does the chapter's scaling give it, and what would a raw loss of 34\% become?
\end{exercise}

\begin{exercise}[$\star$]\label{exo:s1:momentum:2}
After a 50\% loss, what return recovers it? And after a 50\% loss followed by a 100\% gain, where is the strategy?
\end{exercise}

\begin{exercise}[$\star$]\label{exo:s1:momentum:3}
The decile book trades 5.7 times its gross a year one-way. What does ten basis points per unit traded cost it a year?
\end{exercise}

\begin{exercise}[$\star\star$]\label{exo:s1:momentum:4}
Explain why a momentum book is short high-beta stocks after a long market decline, and what happens when the market rebounds.
\end{exercise}

\begin{exercise}[$\star\star$]\label{exo:s1:momentum:5}
Why does momentum skip the most recent month?
\end{exercise}

\begin{exercise}[$\star\star$]\label{exo:s1:momentum:6}
The scaled factor's worst month (September 1939) came with a weight of 0.92. What does that say about what volatility scaling can and cannot do?
\end{exercise}

\begin{exercise}[$\star\star\star$]\label{exo:s1:momentum:7}
\emph{Coding.} Run \texttt{scaled('residual')}. Compare the Sharpe ratios with the total-return book's, and explain why scaling helps the French factor
more than the synthetic books.
\end{exercise}

\begin{exercise}[$\star\star\star$]\label{exo:s1:momentum:8}
\emph{Find the flaw.} ``Momentum has earned 7\% a year for a century with a Sharpe ratio of 0.45; at three times leverage it is a 21\% strategy.''
\end{exercise}

\section{Problem: Spring 2009}

\begin{problem}[{Weekend problem --- the crash and its cure}]\label{pb:s1:momentum:1}
The French momentum factor, raw and scaled, and three momentum books on \texttt{firm.synthmkt}.

\medskip\noindent\textbf{Part I --- The premium.}
\begin{enumerate}
\item Define \omterm{def:s1:momentum:cross}{cross-sectional momentum} and the 12--1 signal.
\item What did Jegadeesh and Titman find?
\item Give the French factor's return, volatility, Sharpe ratio and skewness since 1927.
\item What are its worst month and its maximum drawdown?
\end{enumerate}

\medskip\noindent\textbf{Part II --- The variants.}
\begin{enumerate}[resume]
\item Define residual and \omterm{def:s1:momentum:residual}{industry momentum} and give their published results.
\item Give the three synthetic books' ICs and Sharpe ratios.
\item Why does \omterm{def:s1:momentum:residual}{residual momentum} win here?
\item Why does \omterm{def:s1:momentum:residual}{industry momentum} earn nothing here?
\end{enumerate}

\medskip\noindent\textbf{Part III --- Crashes.}
\begin{enumerate}[resume]
\item Define a \omterm{def:s1:momentum:crash}{momentum crash} and describe when crashes happen.
\item Explain the mechanism in terms of beta.
\item Which were the ten worst months?
\item Why does the synthetic market not reproduce the mechanism?
\end{enumerate}

\medskip\noindent\textbf{Part IV --- The verdict.}
\begin{enumerate}[resume]
\item State the \emph{named result}: the momentum factor's drawdown in 2009 in the French data and that of the volatility-scaled version.
\item What does scaling do to the long-run Sharpe ratio and skewness?
\item When does scaling fail?
\item What does the scaled version cost in leverage?
\item What did Gao and co-authors find within the day?
\item What would you monitor on a live momentum book?
\item How much do momentum's costs matter compared with chapter~2's?
\item In one sentence: what is momentum's risk?
\end{enumerate}
\end{problem}

\section{Interview questions}

\begin{interviewq}[$\star$ \iqroles{researcher, trader}]\label{iq:s1:momentum:1}
What is momentum, and why skip the last month?
\end{interviewq}

\begin{interviewq}[$\star\star$ \iqroles{researcher}]\label{iq:s1:momentum:2}
Why does \omterm{def:s1:momentum:crash}{momentum crash}, and when?
\end{interviewq}

\begin{interviewq}[$\star\star$ \iqroles{researcher, risk}]\label{iq:s1:momentum:3}
How would you manage momentum's crash risk?
\end{interviewq}

\begin{interviewq}[$\star\star$ \iqroles{researcher}]\label{iq:s1:momentum:4}
\omterm{def:s1:momentum:residual}{Residual momentum} versus total-return momentum: what is the difference, and which would you trade?
\end{interviewq}

\begin{interviewq}[$\star\star$ \iqroles{risk}]\label{iq:s1:momentum:5}
Your momentum book's beta was 0.1 a month ago and is $-0.4$ today. What happened, and what do you do?
\end{interviewq}

\begin{interviewq}[$\star\star\star$ \iqroles{researcher}]\label{iq:s1:momentum:6}
Show that scaling a strategy by the inverse of a forecast of its volatility raises its Sharpe ratio when volatility is predictable and unrelated to the mean.
When can it lower it?
\end{interviewq}
